\documentclass[12pt]{article}

\usepackage{fullpage}
\usepackage[utf8]{inputenc}

\begin{document}

TP1, ROP631, été 2026, Prof. Toby HOCKING, à rendre le 5 juin.

\section{exercices 1.4.2 et 2.3.5}
 Consid\'erez une hyperbole de $\Re^2$, $y=1/x$. Soit un point
 $P={P_x\choose P_y}$, on veut trouver le point de l'hyperbole le plus
 proche de $P$. Comme les points de l'hyperbole sont de la forme
 ${x\choose 1/x}$, on peut ramener l'\'etude \`a la minimisation selon
 la seule variable $x$ de la fonction $d(x) = (x-P_x)^2 + ({1\over x}-P_y)^2$.
 \begin{enumerate}
 \item \'Ecrivez les conditions d'optimalit\'e (d'ordre un et deux) de ce probl\`eme
   $\min_{x\in\Re}d(x)$.
 \item Pour $P=(2,3)^t$, utilisez Julia pour obtenir la solution.
 \item Si le point $P$ est sur la droite $y=x$, on devine qu'il y a
   2 candidats solution naturels: $1\choose 1$ et $-1\choose -1$. En
   vous limitant aux $x>0$, vérifiez cette
   intuition en utilisant les conditions d'optimalité; vérifiez
   d'abord que $x=1$ est un point stationnaire pour
   $d(x)$. Vérifiez ensuite que
   l'intuition n'est pas vraie si $P$ est assez loin de l'origine;
   quantifiez ce ``loin'' de l'origine en utilisant les conditions
   d'optimalité d'ordre 2.)
 \end{enumerate}

\section{exercice 2.4.5}

\begin{enumerate}
\item \`A partir de $x=1$, appliquez l'algorithme {\tt
    trouve\_intervalle} pour le point $P={1\choose 2}$; utilisez un
  paramètre $\delta=0.5$. \label{iii}
\item \`A partir de l'intervalle trouv\'e en a), effectuez
  une \'etape de l'algorithme {\tt bissec\-tion} (toujours pour le
  point $P={1\choose 2}$). \label{iv}
\end{enumerate}

\section{exercice 2.5.6}

\`A partir de chacune des extrémités de la sortie de \verb|trouve_intervalle|, effectuez une it\'eration de l'algorithme {\tt Newton-Raphson}(toujours pour le point $P={1\choose 2}$).

\section{Implantation d'algorithme}

Il s'agit de produire
du code qui implémente l'algorithme de région de confiance (algorithme 2.5, page 90 des
notes). Pour une fonction donnée $h(t)$ de classe $C^2$, à partir d'un point de départ arbitraire
$t_0$, l'algorithme calcule un minimum local de $h$.
Le fichier \verb|TR_1D_fr.jl| contient une enveloppe du code que vous devrez compléter. Le
fichier \verb|exemple.jl| contient le script que vous appelerez, qui charge une de quatre fonctions
fournies en exemple, exécute \verb|TR_1D|. La version \verb|_fr| comporte des commentaires en français.
Tous ces fichiers étaient dans le TD0.
Vous devez :
\begin{enumerate}
\item compléter le code de l'algorithme dans le fichier \verb|TR_1D_fr.jl| ;
\item enlever le \verb|maxiter = 1| du fichier \verb|exemple.jl| pour tester correctement votre solution ;
\item exécuter votre code sur les quatre exemples fournis ;
\item valider les résultats sur les quatre exemples fournis ;
\item justifier que l'ordre de convergence asymptotique est quadratique. Astuce : lire section 2.6. Pour chaque itération $x_k$, calculer l’erreur $\epsilon_k=x_k-\bar x$, où $\bar x$ est la solution trouvée. Ensuite calculez $|\epsilon_{k+1}|/|\epsilon_k|^p$ pour $p=1$ (linéare) et $p=2$ (quadratique).
\end{enumerate}

\end{document}
